% opencmdb — User Manual. Build: see ../Makefile (LuaLaTeX; TEXINPUTS -> ../common).
\documentclass[11pt,oneside,listof=totoc,bibliography=totoc]{scrreprt}
\usepackage{opencmdb-manual}
\setmanualname{User Manual}

\begin{document}
\makemanualtitle{User Manual}{Understanding and operating opencmdb day to day}

\tableofcontents

\chapter{Introduction}
\productname{} is a self-hosted, single-binary engine that helps you keep an accurate,
living picture of your network. Its core idea is a comparison:

\begin{itemize}
  \item the \textbf{observed} state — what opencmdb \emph{discovers} automatically from
        your network (devices, IP addresses, switch ports, DHCP leases, \ldots);
  \item the \textbf{declared} state — what \emph{you} have documented as intended.
\end{itemize}

The difference between the two — the \textbf{gap} — is what opencmdb surfaces for you to
review. Observed data is never allowed to overwrite what you declared: the two are kept as
distinct, linked records.

\begin{note}
This manual is written for the person who \emph{uses} opencmdb — reviewing discoveries,
documenting devices, and managing IP space. Installation, configuration, security, and
backups live in the companion \textbf{Administrator Manual}.
\end{note}

\section{Who it is for}
Advanced home labs and small organisations without a dedicated IT team, typically running
on a Synology NAS. The reference scale is on the order of a few hundred hosts across a few
dozen subnets.

\section{What it is not}
opencmdb does not reconfigure your network and does not require elevated network
privileges to observe it. It documents and reconciles; it does not control.

\chapter{Core concepts}
\section{Declared vs.\ observed}
Every entity in opencmdb has (up to) two faces: a \emph{declared} record you author, and an
\emph{observed} record discovery produces. They are linked, never merged. Reconciliation
\emph{links} an observation to the declaration it matches; it never rewrites your intent.

\section{Device identity}
A device is identified by a composite identity — not by its MAC address alone — so it stays
recognisable across a changing IP or MAC. When the evidence is ambiguous, opencmdb
\textbf{abstains} rather than guessing or merging: an unresolved match is a first-class,
recorded outcome, presented with its supporting evidence.

When the engine will not say whether several network interfaces are one machine --- for
instance two cards answering to the same name --- the triage inbox shows \emph{one}
\textbf{Ambiguous} line for them. Its pane lists each candidate's hardware address, its
addresses and names \emph{as the network shows them now}, and why the engine asks rather than
decides. A shared name is weak evidence on its own, which is why it is a question and never a
merge. The question follows the network: it leaves the inbox if the names stop matching, and
returns when they match again. The addresses concerned keep their own \emph{Add} control ---
where documenting is enabled on the instance --- and say that they belong to an open question;
adding one creates a record for that address alone and does not answer the question.

\begin{planned}
The identity engine groups nothing yet: it recognises a network interface by its hardware
address, and when several interfaces may be one machine it \emph{asks} (below) rather than
grouping them. Composite identity across a changing address is the designed model.
\end{planned}

\section{Sources and liveness}
Each discovery source reports on two independent axes: \emph{liveness}
(\texttt{live} or \texttt{blind}, with a named cause) and \emph{capability} (what it can
currently see). When a source goes blind, opencmdb freezes its last-known state and
suppresses observation-derived alerts for that scope — it never fabricates a divergence from
a source that cannot see.

\begin{planned}
The \emph{sources} screen reports which fact kinds a source can and cannot see, measured at
runtime. The \emph{liveness} axis carries no colour yet: the product can establish that an
observation arrived, and cannot establish that a source has gone blind --- deriving that from
silence would invent an incident. The incident axis arrives with Epic~13.
\end{planned}

\section{Triage gestures}
When you review a discovery, you choose one canonical gesture:
\begin{itemize}
  \item \textbf{Document} — close the gap by adopting the observed value into your
        declaration, field by field. You document a \emph{value}.
  \item \textbf{Accept-gap} — record that you have \emph{seen} the difference but are not
        deciding yet. The gap stays open and keeps counting; a note is required. You accept a
        \emph{gap}.
  \item \textbf{Exclude} — this item is out of scope and should stop appearing.
  \item \textbf{Attach}, \textbf{create}, \textbf{snooze} — link to an existing record,
        create a new declaration, or defer.
  \item \textbf{Resolve} — answer an \emph{ambiguity}: say whether the interfaces it names
        are one machine or several. See \emph{Answering an ambiguity} below.
\end{itemize}

\begin{note}
opencmdb remembers your triage decisions (including exclusions), so a resolved item does not
reappear on the next scan.
\end{note}

\begin{planned}
The triage screen shows the gap gestures today, each labelled \emph{Not yet}: they do not
act yet. They are visible so that the shape of the work is legible before the gestures are
built. Epic~7 makes them live.
\end{planned}

\subsection{Answering an ambiguity}
An \emph{Ambiguous} line asks whether several network interfaces --- two cards answering to
one name, for example --- are one machine or several. Its pane shows each candidate as the
network shows it now, says why the engine asks rather than decides, and carries
\textbf{Resolve} as two answers:
\begin{itemize}
  \item \textbf{The same machine} --- the candidates are one box.
  \item \textbf{Distinct machines} --- none of them is the other.
\end{itemize}
Either answer is kept as \emph{yours}: the engine will not ask that question again and will
never overwrite your answer. The question leaves the queue at once, and the addresses it
concerned offer \textbf{Add} again. While the question is open, those addresses show a link
to it \emph{instead of} \textbf{Add}, so that you answer before you record a machine twice.
The reach section at the bottom of the screen counts the questions you have answered.

If some of the interfaces were already answered as \emph{distinct machines}, only
\textbf{Distinct machines} is offered: saying \emph{the same machine} now would contradict your
earlier answer, and the pane says so.

\begin{note}
The answer is dated at the moment the page showed you, not at the moment you pressed. If the
network changed the question in between, or its machines changed, the product refuses the
answer and asks you to reload and look again --- nothing is written.
\end{note}

\begin{planned}
Two limits of this version. An answer cannot yet be changed or undone from the screen, and the
count says \emph{how many} questions you answered, not \emph{which}. And the inventory still
lists each address you added on its own: a machine you called one appears once only when the
device grouping arrives --- so after answering \emph{the same machine}, pressing \textbf{Add} on
each of its addresses records it twice, which is exactly what answering first was meant to avoid. A question whose machines are not answering on the network right now
offers no answer, and says so; the answers return with the next sweep that sees them.
\end{planned}

\chapter{Getting started}
\begin{planned}
A guided first-run walkthrough will appear here once the application is runnable. For now,
deployment is covered in the Administrator Manual (Docker on Synology is the priority
target).
\end{planned}

\section{First light}
On first launch you will connect at least one discovery source and let opencmdb build its
initial observed picture. Nothing is alarming on day one: undocumented discoveries simply
wait for you in the triage inbox.

\chapter{Discovering your network}
\section{The UniFi source}
If you run a UniFi controller, opencmdb can discover devices, IP addresses, switch ports,
SSIDs, VLANs, and DHCP leases from it using only an API key — no elevated network access.

\begin{planned}
There is no UniFi connector yet. The only source that ships is the ARP/ping scanner below,
which sees an IPv4 address and a round-trip time and nothing else --- not even a hardware
address. The \emph{sources} screen names the five fact kinds it cannot see.
\end{planned}

\section{The ARP/ping scanner}
For everything else, declare the subnets to scan and opencmdb will find active devices on
them, enriching with hostnames where available.

\begin{planned}
Both sources are designed and specified; the ARP/ping scanner is the first source to be
implemented, followed by the full UniFi connector.
\end{planned}

\chapter{The triage inbox}
The inbox is where undocumented or changed discoveries wait for your decision. It shows, in
priority order, what changed since your last visit — new devices and genuine conflicts
first, routine churn last.

\section{Documenting a discovery}
\begin{itemize}
  \item \textbf{Document-all} adopts a brand-new discovery in a single action.
  \item \textbf{Document-field} lets you adopt specific fields of a device whose observed and
        declared values have diverged since you last added it. The observed record itself is
        never modified.
\end{itemize}

\textbf{The interface calls this gesture \emph{Add} / « Ajouter »}, not \emph{Document}. The
identifier stays \texttt{document} in this manual, in the API and in the code --- it is the one
concept whose label and identifier differ, because the engine can only ever ADD a declared value
and never replace one.

\textbf{What it creates, and what it does not.} The gesture copies what the connector observed
--- today an address, a hostname when reverse DNS answers, and a hardware address. \textbf{You
cannot yet write a field the network did not show you}: there is no form for a name, a role or an
owner. So a record is what the sighting carried, and nothing more.

\begin{planned}
\emph{Document-all} has been reachable from the triage queue since \texttt{v0.3.0}. What is not
built is \emph{document-field}, the per-field gesture, and with it any way to enter a value of
your own; Epic~7 gives \emph{it} a control the operator can press.
\end{planned}

\section{Working in bulk}
You can triage many discoveries at once when they share a decision, keeping steady-state
review fast.

\begin{planned}
The inbox, the document/accept-gap/exclude gestures, and bulk triage are specified in the UX
design and will be implemented over several milestones.
\end{planned}

\chapter{IP address management}
opencmdb manages subnets, VLANs, and DHCP ranges, and shows per-subnet occupancy — used,
free, declared, and observed. It can find a free address within a subnet and detect
conflicts (the same IP on two MACs, or a statically declared IP inside a DHCP range),
identifying the devices involved. IPv6 addresses can be documented (observation-only at
this stage).

\section{Defining the addressing plan}
The \emph{IPAM} screen draws the plan you define, one subnet at a time. On a fresh install the
plan is empty; the screen says so and links to the form that fills it. From the
\emph{Change the plan} panel you can:
\begin{itemize}
  \item define a \textbf{subnet}, in CIDR notation (for example \texttt{192.0.2.0/24}), with a
        label and, optionally, a \textbf{VLAN} --- a number from 1 to 4094, or nothing at all.
        The same address space may be declared once per VLAN, so \texttt{192.0.2.0/24} on VLAN~10
        and the same CIDR on VLAN~20 are two entries of the plan; declaring it twice in one VLAN is
        refused;
  \item inside the selected subnet, define a \textbf{range}: a first and a last address, a label,
        and what the range is meant for --- \texttt{static}, \texttt{dhcp-pool},
        \texttt{reserved} or \texttt{infrastructure};
  \item inside the selected subnet, define an individual \textbf{address}, with a label.
\end{itemize}
A refused write says which rule it broke: a range that overlaps another or leaves its subnet, a
subnet written with host bits set, a label that is empty or longer than 120 characters, a record
the plan already holds. The forms need JavaScript. If a write takes too long, the screen says it
may or may not have been saved: reload the plan before trying again.

Beside the grid, two lists show what the selected subnet already holds --- its declared ranges and
its defined addresses --- and each entry carries two controls, \emph{Correct} and \emph{Remove}.
Correcting keeps the record: the form arrives holding its current values, so you change what is
wrong rather than retyping the rest. The rules that governed the definition still apply --- a range
may be widened onto ground it already occupies; it may not be widened onto a neighbour's --- and
correcting a range carries one rule that defining it does not: a range may not be moved or narrowed
off an address you defined inside it. Widening is therefore always allowed, and so is changing the
policy or the label; what is refused is a correction that would leave one of your own addresses
outside the range that explains it. These lists are
shown whether or not the grid is drawn, so a subnet too large to draw still offers the controls that
correct it.

Removing is refused when the removal would leave something behind: a range that still holds
addresses you defined inside it says so by name, and a subnet that still holds ranges or addresses
says that it holds them --- not that it does not exist. Before any of the three removals fires,
reaching its control tells you what that removal would change: for a range, how many addresses you
defined inside it still hold it back --- which is the refusal above, announced before you meet it
rather than after --- how many addresses the network shows inside it are explained by it today and
would become findings once it is gone, and whether it is the only \texttt{static} range, so that
removing it would leave nothing to offer here; for a defined address, whether the network still
shows it; and for a subnet, how many records are still holding it back. As with the other two
warnings, \emph{the warning} refuses nothing: it says what would change and never stops the gesture
itself. Whether the removal is then refused is the paragraph above --- a range still holding your
addresses, and a subnet still holding anything, are both refused until what they hold is gone.

Two of the forms tell you what they know \emph{before} you write. As you type an address, the form
says whether the network has been seen on it (and with which hardware addresses, and when), whether
the inventory documents it, whether the plan already defines it, and whether it falls inside a
\texttt{dhcp-pool} range. As you type a range meant for \texttt{static} use, the form says how many
addresses already seen on the network would fall inside it. Neither warning refuses anything: writing
down the machine that is already there is the ordinary case, and the last line of each says so. If
the check cannot reach the store it says that too, rather than leaving you the previous answer.

\section{What the network shows against the plan}
Under the grid, \emph{What the network shows} lists every address of the selected subnet the
scanner has seen that the plan has something to say about:
\begin{itemize}
  \item \textbf{gap} --- the plan would offer the address as free, and a machine is using it;
  \item \textbf{undeclared} --- the address is in use, has no entry in the plan, and is not one
        the plan would offer: inside a \texttt{reserved} or \texttt{infrastructure} range, on the
        network or broadcast address, or where no range says anything;
  \item \textbf{Address conflict} --- two hardware addresses were seen on one address inside a
        \texttt{static} or \texttt{reserved} range. A replaced network card and a duplicate look the
        same from the network, so the entry names both hardware addresses and when each was last
        seen, and leaves the judgement to you.
\end{itemize}
Nothing inside a \texttt{dhcp-pool} range is listed: there, the machine using an address changes by
design. Each entry gives, for every hardware address seen, the date it was last seen --- a date, not
an age --- and links to its question in the triage screen when there is one. An address documented
in the inventory has no triage question, and the entry says so.

A sighting protects an address \emph{until you release it}: the product cannot tell a machine that
left from one that is asleep, so an address seen once stays listed. The list also names addresses
seen outside every subnet of the plan, and addresses you defined inside a \texttt{dhcp-pool} range,
which the pool may hand to another machine.

\section{Releasing an address}
Each \textbf{gap} or \textbf{undeclared} entry of the selected subnet carries a \emph{Release}
control, named after its address. Releasing tells the plan to forget what the network had shown on
that address \emph{when your page was drawn} --- the sightings listed in the entry, and nothing that
arrived after it. The entry leaves the list, the page says the address was released, and inside a
\texttt{static} range the address can be offered again --- unless the inventory documents it, since
a documented address is never offered. Nothing the scanner recorded is deleted --- the release is
kept as a record of its own, beside the sightings.

A release lapses by itself when the network answers again. A machine seen on the address after the
last sighting your page showed holds the address once more, and the entry comes back as an ordinary
finding --- an address something answers on is never offered. A network card replaced after the
release counts only from the moment it was seen, so it is not reported as a conflict with the card it
replaced.

An address you defined in the plan cannot be released: the plan never offers it anyway, and the
refusal says to remove the address instead. Defining an address you had released cancels the
release. A release cannot be undone from the screen; if you released a machine that was only asleep,
its address may be offered until it is seen again --- the Administrator Manual says how to remove a
release by hand.

\section{The next address the plan can offer}
The panel beside the grid proposes the lowest address that nothing the product knows of speaks for:
inside a \texttt{static} range, not the network or broadcast address, not defined in the plan,
\emph{never seen on the network}, and not documented in the inventory. It is not a promise that the
address is free: the product can only exclude what it has been told or has seen, so a machine that
was asleep during every scan, or that lives outside the range you configured the scanner to sweep,
is invisible to it and keeps its address all the same. When several ranges cover an address ---
overlapping ranges, or a subnet nested in another --- the most protective one decides: a
\texttt{reserved} range over a \texttt{static} one takes the address out of the offer. A subnet too
large to draw proposes nothing and still lists what the network shows. The occupancy line counts as
free only what the offer would propose.

\section{What the VLAN does, and what it does not}
The VLAN is part of what \emph{you declare}. The scanner does not read 802.1Q tags, and nothing in
the product relates a VLAN you wrote down to a broadcast domain it observed --- so the audit ignores
it entirely. Ranges, findings and the free-address offer are computed across the whole plan, and two
subnets sharing one address space share one audit: a \texttt{dhcp-pool} you declare in one of them
changes what the other reports. The screen carries that sentence under the subnet selector, so you
meet it where it matters rather than here. What the VLAN buys you today is the ability to write down
a plan that reuses one address space in several segments, and a selector that tells those segments
apart.

\section{Correcting a subnet}
A subnet's own \emph{Correct} control, at the foot of the panel beside the grid, changes its label
and its VLAN --- and nothing else. There is deliberately no field for its address space: moving a
subnet's CIDR would leave the ranges and addresses hanging from it outside their own parent, so the
way to change an address space is to define the subnet you want and remove the one you do not.
Clearing the VLAN field means \emph{no VLAN}; a subnet cannot be moved onto a VLAN another entry of
the same address space already holds.

\section{IPv6}
opencmdb's plan holds IPv6 subnets, ranges and addresses. You write them as you write IPv4 ones ---
\texttt{2001:db8::/64}, a first and a last address, a label --- and the selector lists them beside
your IPv4 subnets.

\textbf{What an IPv6 subnet's page shows, and why it shows less.} The scanner does not reach IPv6:
it reads hardware addresses off the IPv4 neighbour table, and there is no IPv6 equivalent in this
product. So for an IPv6 subnet there is nothing to compare your plan against, and the page says so
rather than drawing an empty audit. You get the ranges and addresses you declared; you do not get a
grid, an occupancy count, a free-address offer, or a list of findings.

That is deliberate. An empty findings list reads as \emph{nothing is wrong}, where the truth is
\emph{this subnet was not compared against the network} --- and the second sentence is the one you
need. For the same reason, releasing an IPv6 address is refused: a release makes the plan forget what
the network showed on an address, and the product has never seen anything on yours.

What you do still get is every check that needs no observation. If you define an address inside a
range you declared a DHCP pool, the warning appears on an IPv6 subnet exactly as it does on an IPv4
one: that is a contradiction inside your plan, and reading the network would not help settle it.

\textbf{The list of addresses seen outside every subnet of your plan still appears} on an IPv6
subnet's page. It belongs to the plan rather than to the subnet you are looking at, so hiding it
would cost you a true finding for having selected a tab. It carries its own heading saying so.

\textbf{Write an IPv4 space as IPv4.} \texttt{::ffff:192.0.2.0/120} is a legal way to write
\texttt{192.0.2.0/24} in IPv6 notation, and opencmdb refuses it: held as an IPv6 subnet it would
tell you the product cannot observe a range the scanner sweeps every five minutes. Enter it as IPv4
and you get the grid, the findings and the offer.

\chapter{Applications and hosted software}
Beyond the network layer, you can record software instances (name, version, listening ports)
hosted on a device, and group them into applications with an owner and a criticality. This
is the ``light CMDB'' layer: enough to answer \emph{what runs where, and who owns it}.

\begin{planned}
The applications layer is planned for a later milestone, after the core reconciliation loop.
\end{planned}

\chapter{Alerts and notifications}
opencmdb can notify you of the things that matter — new conflicts, a source going blind, a
capability downgrade — without drowning you in routine change. Alerts derived from
observation are suppressed for any source that is currently blind.

\begin{planned}
The alerts screen ships as an \emph{example}: the three alert kinds it lists are the ones the
product will raise, and no alert is raised or delivered today. Epic~13 builds them.
\end{planned}

\chapter{A bilingual interface}
The user interface is available in English and French; documentation is English-only. Set
\texttt{OPENCMDB\_LOCALE} to \texttt{en} or \texttt{fr} (a region suffix such as
\texttt{fr-CH} is accepted); an unrecognised value refuses to start, naming the variable,
rather than falling back to English in silence. The interface is \textbf{light}, on the
reference design's palette --- \emph{this manual said a dark theme was the default until
\texttt{v0.2.0}}. Key views target WCAG~2.1~AA: every screen is measured by axe-core on each
CI run, and the triage queue is reachable by keyboard.

\appendix
\chapter{Glossary}
\begin{description}[style=nextline,leftmargin=0pt]
  \item[Declared record] What you have documented as intended. Never overwritten by
        discovery.
  \item[Observed record] What discovery found. Linked to a declaration, never merged into it.
  \item[Gap] A difference between a declaration and its linked observation — the unit of
        work opencmdb surfaces.
  \item[Document] Close a gap by adopting an observed value into your declaration.
  \item[Accept-gap] Acknowledge a gap without resolving it; it stays open and keeps counting.
  \item[Release] Stop the plan holding an address: what the network had shown on it is forgotten,
        and the address can be offered again. A release is kept as a record, never a
        deletion. Not a software release --- this is the one term with that second meaning.
  \item[Blind source] A source that currently cannot see; its last-known state is frozen and
        its alerts suppressed.
  \item[Abstention] A recorded ``unresolved'' outcome when identity is ambiguous — never a
        guess or a merge.
\end{description}

\end{document}
